% ---------------------------------------------------------------------------
%  Author Summary — required by PLOS journals (Neglected Tropical Diseases,
%  Global Public Health), not by Epidemics or Royal Society Open Science.
%
%  PLOS asks for 150-200 words aimed at a reader outside the field: no
%  statistics, no jargon, and the significance stated plainly. Kept in its own
%  file so the manuscript compiles unchanged for venues that do not want it.
%
%  Word count: 192.
% ---------------------------------------------------------------------------

\section*{Author summary}

Whether dengue outbreaks are driven by weather is a question researchers answer
by fitting a mathematical model of transmission to reported case counts, twice:
once allowing temperature and rainfall to matter, once not, and then asking which
version fits better. Answering it that way requires decisions the model itself
does not settle --- how to describe the randomness in weekly counts, how many
weeks after rain to expect cases, whether to represent mosquitoes explicitly.
Papers rarely say which decisions they made.

We fitted every suitable dengue outbreak in a global surveillance database ---
221 of them, across 33 countries --- under all 144 combinations of six such
decisions. The conclusion changed with those decisions far more often than it
changed from one outbreak to another. Most of that instability was not a matter
of some methods being consistently more permissive: it depended on which method
met which epidemic, which means no agreed standard would remove it.

We show that reporting a result from a handful of alternative analyses instead
of one restores the reliability that a single analysis promises but does not
deliver, at a cost of a few minutes of computing.
